
\section{Evaluation}
\label{04_Results&Evaluation}

To extend the evaluation of SECDA-DSE, we generated and evaluated three accelerator workloads: element-wise vector multiplication, 2D convolution, and matrix transpose. These workloads were intentionally selected to represent different accelerator execution characteristics, including arithmetic intensive computation, convolution-oriented spatial compute, and memory movement dominated execution patterns. The prompts used for generating these accelerator designs are provided in the Appendix. %~\ref{sec:appendix-prompt}.

The SECDA-DSE workflow begins with a natural language accelerator specification provided by the user. This prompt is processed through the LLM Stack, where the RAG module retrieves the relevant SECDA implementation context, architectural templates, and previously collected hardware datapoints from the vectorized knowledge-base. The retrieved context is then combined with CoT prompting to generate a structured accelerator generation prompt enriched with SECDA-specific architectural guidance.

Using this enriched context, SECDA-DSE generates native SystemC accelerator designs together with the required SECDA integration files. In the current implementation, the generated designs are passed through a human-in-the-loop evaluation workflow using the SECDA toolkit. The generated accelerators are first validated using HLS, and the designs that successfully pass HLS are then evaluated through downstream logic synthesis and FPGA execution flows.

The evaluation was performed using Vivado HLS 2019.2~\cite{vivadohls2019_2} targeting a Xilinx Zynq-7000 FPGA platform (\texttt{xc7z020-clg400-1}). The LLM Stack used TinyLlama~\cite{zhang2024tinyllama}, a lightweight 1.1B parameter model deployed locally through Ollama~\cite{ollama}. Initially, the LLM was only fine-tuned using hardware datapoints generated from matrix addition and matrix multiplication accelerator workloads.

Table~\ref{tab:workload_comparison} summarizes the hardware execution and FPGA utilization characteristics for the generated workloads. The primary objective of this evaluation was to validate SECDA-DSE across multiple accelerator workloads through the complete hardware generation workflow, including HLS, logic synthesis, and FPGA execution. %To verify the correctness of the generated designs, the output obtained from the FPGA is compared element by element with the reference result computed directly by the CPU using loops.
The correctness of the generated accelerator designs is validated by performing an element-wise comparison between the outputs produced by the FPGA and reference results generated by a CPU-based software implementation.
At this stage, the focus is on end-to-end hardware execution rather than producing fully optimized accelerator implementations. As discussed in more detail in Section~\ref{Limitation and future work}, subsequent iterations of SECDA-DSE will focus on performance-optimized accelerator designs.


\begin{scriptsize}
\begin{table}[t]
\centering
\caption{FPGA execution and resource utilization comparison for generated accelerator designs.}
\label{tab:workload_comparison}
\setlength{\tabcolsep}{3pt}
\renewcommand{\arraystretch}{1.1}
\begin{tabular}{|l|c|c|c|}
\hline
\textbf{Metric} & \textbf{2D Conv} & \textbf{VMUL} & \textbf{Transpose} \\
\hline
\hline
Validation & PASSED & PASSED & PASSED \\
\hline
Latency (ms) & 163 & 135 & 238 \\
\hline
HWC cycles (1/2/3) & 1251/76/1250 & 52/26/51 & 393/311/387 \\
\hline
\hline
DMA recv size (bytes) & 256 & 64 & 512 \\
\hline
DMA send size (bytes) & 268 & 136 & 524 \\
\hline
DMA recv speed (MB/s) & 17.07 & 4.57 & 13.47 \\
\hline
DMA send speed (MB/s) & 38.29 & 19.43 & 30.82 \\
\hline
DMA recv wait (ms) & 15 & 14 & 38 \\
\hline
DMA send wait (ms) & 7 & 7 & 17 \\
\hline
\hline
BRAM utilization (\%) & 2.50 & 3.57 & 2.86 \\
\hline
DSP utilization (\%) & 1.36 & 21.82 & 5.45 \\
\hline
LUT utilization (\%) & 6.64 & 8.23 & 8.85 \\
\hline
FF utilization (\%) & 4.30 & 6.02 & 5.84 \\
\hline
% WNS & -1.57 & -2.44 & -2.39 \\
% \hline
% TNS & -161.69 & -1258.88 & -493.64 \\
% \hline
\end{tabular}
\end{table}
\end{scriptsize}


Despite being generated from natural language specifications, the produced accelerator designs were executable without requiring manual modification of the generated accelerator logic. The generated workloads exhibited distinct hardware characteristics depending on the computational structure of the workload. For example, vector multiplication workload emphasized arithmetic execution and parallel compute behavior, resulting in higher DSP utilization (\(21.82\%\)) compared to the convolution (\(1.36\%\)) and transpose (\(5.45\%\)) accelerators. While the transpose workload demonstrated a more memory movement dominated execution pattern with increased DMA transfer sizes and higher FPGA logic utilization (\(8.85\%\) LUT usage).

Using hardware counter (HWC) cycles and DMA profiling metrics, SECDA-DSE identifies bottlenecks within generated accelerator designs and uses this information during iterative refinement. HWC1 measures data loading wait cycles, HWC2 captures computation time, and HWC3 records write-back cycles. 
The 2D convolution accelerator demonstrated substantially higher cycle counts (\(1251/76/1250\)) compared to the vector multiplication (\(52/26/51\)) due to the increased buffering and data reuse complexity, while the transpose accelerator showed intermediate execution complexity (\(393/311/387\)) because of runtime matrix reorganization overheads. 
Similarly, the transpose workload required the largest DMA transfers (\(512\) receive bytes and \(524\) send bytes), reflecting its memory-centric execution pattern, whereas vector multiplication emphasized arithmetic throughput with smaller transfers. These profiling metrics are stored as hardware datapoints and reused by the LLM Stack during subsequent refinement iterations.

SECDA-DSE generated these accelerator designs using a relatively small 1.1B parameter model with only limited initial fine-tuning datapoints. Despite the constrained training setup, the framework successfully generated FPGA executable accelerator designs across multiple workload types while maintaining SECDA compatibility and valid FPGA execution flows.

The three generated accelerator workloads also exhibited different convergence behaviors during iterative refinement. The vector multiplication required four iterative refinement cycles before producing a valid FPGA executable design. 
In contrast, the 2D convolution accelerator successfully generated a valid design in a single generation cycle. The transpose accelerator required nine refinement cycles overall, where the first two iterations produced an HLS synthesizable design while an additional seven iterations were required to successfully complete downstream logic synthesis and FPGA execution. Throughout these iterations, human intervention was restricted to operational tasks required by the SECDA validation workflow, such as transferring generated source files, verifying execution paths and running evaluation scripts to collect hardware metrics. The generated accelerator designs were not manually modified at any stage.


% During these iterations, SECDA-DSE continuously refined generated accelerator configurations based on synthesis and execution feedback collected.

These results provide early evidence that combining structured DSE, retrieval-grounded reasoning, and iterative hardware evaluation can enable adaptive accelerator generation workflows capable of handling diverse workload characteristics while reducing the manual effort traditionally required during FPGA accelerator development.